\begin{abstract}
A program that orchestrates calls to a large language model (LLM) can leak a
secret without the secret reaching a prompt: if the secret decides which
requests are sent, whoever sees the network traffic or the bill can learn it.
Resource-aware noninterference handles such channels by charging each operation
a cost determined by input sizes. We show that this cannot work for LLM calls. Any analysis that sees prompt text only through a size measure is
either unsound for providers whose answers depend on content, or rejects a
secret-guarded branch whose two arms are identical. The theorem is mechanised in
Coq and explains why two size-based rules of our own failed. We compare
requests by content instead, resolving each secret branch by the outcomes its
secret can realise. If every resolution sends the same requests, the workflow is
noninterfering for every provider, tokenizer and observer of its transcript; if
$k$ request classes remain, it leaks at most $\log_2 k$ bits of min-capacity.
Because tokenization is not subadditive, we compute worst-case token bounds in
bytes and convert them through contracts measured on fourteen tokenizers. Both
analyses are implemented in OrchLang, a small workflow language, and evaluated on
synthetic suites and on thirty workflows ported from
LangGraph, the Anthropic cookbook and AgentDojo under a protocol fixed before
porting. No certified bound was exceeded in 6,000 runs of the ported workflows.
None of them has a secret, so the relational analysis has so far been exercised
only on synthetic workflows.
\end{abstract}
